\bookchapter{Capítulo 64: Armazenamento de dados e requisições de rede}

\chapterintro{Um app que só “mostra” está morto; app de verdade é o que guarda dados e se conecta à rede. Neste capítulo você vai conhecer o “trio” da camada de dados do iOS: UserDefaults (para pequenas configurações), SwiftData (para dados estruturados) e URLSession (para conversar com servidores), além do Codable para interpretar JSON e do moderno async/await. Com este capítulo dominado, você já escreve apps de verdade, que lembram os hábitos do usuário e atualizam conteúdo em tempo real.}

\booksection{Como escolher a solução de armazenamento: comece por esta tabela}

\begin{center}
\begin{tabularx}{\linewidth}{|X|X|X|X|}
\hline
\textbf{Solução} & \textbf{O que guarda bem} & \textbf{Características} & \textbf{Onde usamos no livro} \tabularnewline
\hline
UserDefaults & Configurações esparsas (interruptores, preferências) & Pares chave-valor, simples e direto & Ajustes do Capítulo 65, cidades favoritas do Capítulo 66 \tabularnewline
SwiftData & Dados estruturados (lançamentos, registros) & Encapsulamento do banco oficial; declaração em uma linha com @Model & Lançamentos do Capítulo 65 \tabularnewline
URLSession & (não armazena, busca na rede) Dados do servidor & Framework de rede nativo, sem dependência de terceiros & Dados de clima do Capítulo 66 \tabularnewline
\hline
\end{tabularx}
\end{center}

\begin{center}\small\color{gray}A divisão de trabalho do trio de dados\end{center}

Grave uma máxima: “configuração pequena vai para o UserDefaults, dados em lote vão para o SwiftData, e dados de fora chegam pelo URLSession”. O Core Data dos tutoriais antigos ainda funciona (o SwiftData é o seu casaco moderno), e SQLite escrito à mão você praticamente não vai tocar — aprenda o novo e não trilhe o caminho velho.

\booksection{UserDefaults: guardando pequenas configurações}

\begin{lstlisting}[language=Swift]
// Escrita: salvar um valor
UserDefaults.standard.set(true, forKey: "isDarkMode")
UserDefaults.standard.set("Pequim", forKey: "favoriteCity")

// Leitura: pegar o valor (se não houver, retorna o padrão)
let isDark = UserDefaults.standard.bool(forKey: "isDarkMode")
let city = UserDefaults.standard.string(forKey: "favoriteCity") ?? "Pequim"

// Exclusão: remover uma chave
UserDefaults.standard.removeObject(forKey: "isDarkMode")
\end{lstlisting}

Lembre que ele só serve para configurações pequenas e esparsas — dezenas ou centenas de registros devem ir para o SwiftData; não trate o UserDefaults como banco de dados (isso deixa o app lento para abrir e dificulta consultas).

\booksection{SwiftData: dando aos dados uma “casa de verdade”}

O SwiftData é a solução oficial de banco de dados desde o iOS 17, e simplifica ao máximo “criar tabelas e fazer CRUD” — com uma única anotação @Model, a classe de dados vira uma tabela do banco:

\begin{lstlisting}[language=Swift]
import SwiftData

// 1. @Model: esta classe vira automaticamente uma "tabela"
@Model
class Record {
    var amount: Double
    var category: String
    var note: String
    var date: Date

    init(amount: Double, category: String, note: String, date: Date) {
        self.amount = amount
        self.category = category
        self.note = note
        self.date = date
    }
}

// 2. Anexar o contêiner na entrada do app
@main
struct MyApp: App {
    var body: some Scene {
        WindowGroup {
            ContentView()
        }
        .modelContainer(for: Record.self)   // banco pronto
    }
}

// 3. Nas views, leia com @Query e escreva com @Environment
struct RecordListView: View {
    @Query(sort: \Record.date, order: .reverse) var records: [Record]
    @Environment(\.modelContext) private var context

    var body: some View {
        List(records) { r in
            Text("\(r.category) \(r.amount) yuan")
        }
    }

    func add(amount: Double, category: String) {
        let r = Record(amount: amount, category: category, note: "", date: Date())
        context.insert(r)     // insere
        try? context.save()   // grava no disco
    }
}
\end{lstlisting}

Esse molde “@Model + modelContainer + @Query” pode ser copiado integralmente para qualquer projeto seu — basta trocar os nomes dos campos. Quando os dados mudam, o @Query atualiza a interface sozinho. É uma delícia.

\booksection{Codable: “traduzindo” o JSON}

A maior parte dos dados dos servidores vem em JSON. O protocolo Codable do Swift faz a conversão “JSON struct” acontecer sozinha — se os nomes dos campos casam, a mágica funciona:

\begin{lstlisting}[language=Swift]
// JSON: {"name":"Pequim","main":{"temp":26.5},"weather":[{"description":"Ensolarado"}]}
struct WeatherResponse: Codable {
    let name: String
    let main: Main
    let weather: [Weather]
}
struct Main: Codable { let temp: Double }
struct Weather: Codable { let description: String }

// String JSON  struct
let decoder = JSONDecoder()
let data = jsonString.data(using: .utf8)!
let result = try decoder.decode(WeatherResponse.self, from: data)
print(result.weather[0].description)   // "Ensolarado"
\end{lstlisting}

Regra de ferro: olhe o JSON real retornado antes de definir os campos. Muitos bugs de “falha na conversão” vêm de nome de campo errado ou tipo chutado errado (o servidor manda a string “26.5” e você declara Double) — imprima o JSON original e compare: dez segundos e você acha o problema.

\booksection{URLSession + async/await: conversando com o servidor}

Para se conectar, o iOS usa o URLSession nativo do sistema; combinado com o async/await do Swift (funções assíncronas), “esperar a rede” fica tão natural quanto “esperar a encomenda chegar”:

\begin{lstlisting}[language=Swift]
struct WeatherService {
    // Funções assíncronas: async marca "preciso esperar um pouco", throws marca "pode falhar"
    func fetchWeather(city: String) async throws -> WeatherResponse {
        let url = URL(string:
            "https://api.example.com/weather?city=\(city)")!
        let (data, _) = try await URLSession.shared.data(from: url)
        return try JSONDecoder().decode(WeatherResponse.self, from: data)
    }
}

// Local da chamada (dentro de .task):
.task {
    do {
        let weather = try await service.fetchWeather(city: "Pequim")
        // atualiza a tela...
    } catch {
        print("Algo deu errado: \(error)")
    }
}
\end{lstlisting}

\textbf{Alerta antiprejuízo: duas regras de ferro da rede}\par

(1) Acesso à rede é sempre com async/await, nunca bloqueando a thread principal — travar a interface é uma experiência de desejar desinstalar. A concorrência estruturada do Swift torna isso natural, mas não esqueça do try await e do do/catch. (2) O ATS proíbe HTTP sem criptografia por padrão: o iOS exige HTTPS. Para testar serviços locais em http://, é preciso abrir uma exceção temporária no Info.plist (NSAppTransportSecurity), mas o produto final deve ser 100\% HTTPS — a revisão da Apple fica de olho nisso.

\begin{tcbpractice}{Mãos à obra: crie um miniapp “frase do dia”}
\booksection{Erros de rede também precisam de “classificação”: não jogue toda falha no colo do usuário}

Sem internet, tempo esgotado, 404, erro 500 do servidor... erros diferentes pedem mensagens diferentes para o usuário:

\begin{lstlisting}[language=Swift]
func friendlyMessage(_ error: Error) -> String {
    if let urlError = error as? URLError {
        switch urlError.code {
        case .notConnectedToInternet: return "A rede caiu — conferir a conexão?"
        case .timedOut:               return "Tempo limite esgotado — tentar de novo depois?"
        case .cannotFindHost:         return "Servidor não encontrado"
        default:                      return "Problema de rede: \(urlError.code)"
        }
    }
    if let d = error as? DecodingError {
        return "Formato de dados inesperado (a API pode ter mudado)"
    }
    return "Algo deu errado: \(error.localizedDescription)"
}
\end{lstlisting}

Mostre ao usuário “linguagem humana” e guarde a “mensagem original” para o log (imprima o error completo). Não jogue no colo do usuário aquele inglês enigmático do error.localizedDescription — o toque de produto mora nesses detalhes.

Procure uma API gratuita de citações (algo como hitokoto), use URLSession + async/await para buscar a “frase do dia”, exiba-a no centro da tela e adicione um botão “Outra frase”. Depois, salve as frases favoritas numa “coleção” com SwiftData. Esse exercício usa UserDefaults/SwiftData/Codable/URLSession todos de uma vez — terminando-o, os capítulos 7 e 8 serão moleza.

\end{tcbpractice}

\begin{tcbsummary}{Resumo do capítulo}
\begin{itemize}
  \item A máxima: configuração pequena no UserDefaults, dados em lote no SwiftData, dados de fora pelo URLSession
  \item Os três golpes do SwiftData: @Model declara a tabela, modelContainer anexa o contêiner, @Query lê e context.insert escreve
  \item O Codable “traduz” o JSON sozinho: se os campos casam, a mágica funciona
  \item URLSession + async/await: conectar é como esperar a encomenda; o do/catch trata as falhas
  \item Duas regras de ferro: não bloqueie a thread principal; produto final exige HTTPS (ATS)
\end{itemize}

\end{tcbsummary}
